\documentclass[../../main.tex]{subfiles}
\begin{document}

\begin{flushright}
  \footnotesize\textit{【自動文字起こし・要確認】}
\end{flushright}

{\bf [解答]}
$Q(x)=0$は重解を持たないと仮定する．$Q(x)=A(x-\alpha)(x-\beta)$（$\alpha\ne\beta$，$A\ne0$）とおく．

$P(x)$は$Q(x)$で割り切れないから，$P(x)$は$(x-\alpha)$，$(x-\beta)$の少なくとも一方を因数に持たない．

\begin{table}[htb]
  \centering
  \begin{tabular}{c|cc}
    因数 & $P(x)$ & $P(x)^2$ \\
    \hline
    $x-\alpha$ & 持たない場合がある & その場合は持たない \\
    $x-\beta$ & 持たない場合がある & その場合は持たない \\
  \end{tabular}
  \caption{$P(x)$と$P(x)^2$の因数の対応}
\end{table}

いずれの場合も，$P(x)$がその因数（例えば$(x-\alpha)$）を持たなければ，$P(x)^2$も$(x-\alpha)$を因数に持たず，$Q(x)$（$(x-\alpha)$を因数に持つ）で割り切れない．

これは，$P(x)^2$が$Q(x)$で割り切れるという仮定に矛盾する．よって$Q(x)$は重解を持つ．$\blacksquare$

\newpage

\end{document}
